A la fin de ce chapitre, je saurai :
\begin{competences}
  \point Définir un signal modulé en amplitude, en fréquence, en phase.
  \point Citer les ordres de grandeur des fréquences utilisées pour les signaux radio AM, FM, la téléphonie mobile.
  \point Interpréter le signal modulé comme le produit d’une porteuse par une modulante.
  \point Décrire le spectre d’un signal modulé.
  \point À partir de l’analyse fréquentielle, justifier la nécessité d’utiliser une opération non linéaire.
  \point Expliquer le principe de la démodulation synchrone.
  \point \faSignLanguage Réaliser une modulation d’amplitude et une démodulation synchrone avec un multiplieur analogique.
\end{competences}
